% Source: https://artofproblemsolving.com/community/c3h2127073_candles_source_expii_solve_set_47_q1
% Stable identity: aops-topic-2127073-q01
You have three candles. Each burns for an exact, known total time: the small one for $1$ hour, the medium one for $3$ hours and the large one for $4$ hours. However, they do not burn at constant rates, so you cannot tell when a given fraction of a candle has burned. Each candle can be lit only once, either now or at the moment another candle burns out, and once lit it burns until it is finished. For example, you can measure $5$ hours from now by burning the $1$-hour candle and lighting the $4$-hour candle as soon as it burns out.

Which of the following is the longest length of time, starting from now, that you \emph{cannot} measure exactly?

$\textbf{(A) }1\text{ hour}\qquad\textbf{(B) }2\text{ hours}\qquad\textbf{(C) }3\text{ hours}\qquad\textbf{(D) }4\text{ hours}\qquad\textbf{(E) }6\text{ hours}\qquad\textbf{(F) }7\text{ hours}$
